\documentclass[12pt,a3paper]{ctexart}

% === 包 ===
\usepackage[top=2cm, bottom=2cm, left=2.5cm, right=2.5cm]{geometry}
\usepackage{graphicx}
\usepackage{float}
\usepackage{caption}
\usepackage{amsmath,amssymb}
\usepackage{hyperref}
\usepackage{fancyhdr}
\usepackage{tikz}
\usetikzlibrary{positioning, shapes.geometric}
\usepackage{booktabs}
\usepackage{longtable}
\usepackage{enumitem}
\usepackage{xcolor}
\usepackage{listings}
\usepackage{multirow}
\usepackage{makecell}
\usepackage{ulem}

\lstset{
    basicstyle=\small\ttfamily,
    breaklines=true,
    frame=single,
    backgroundcolor=\color{gray!3},
    numbers=left,
    numberstyle=\tiny\color{gray},
    commentstyle=\color{green!40!black},
    stringstyle=\color{red!60!black},
    showstringspaces=false,
    tabsize=2
}

\pagestyle{fancy}
\fancyhf{}
\fancyhead[L]{船舶平台综合环境应力剖面分析系统}
\fancyhead[R]{详细设计说明书 V2.0}
\fancyfoot[C]{\thepage}
\renewcommand{\headrulewidth}{0.5pt}

\title{\textbf{船舶平台综合环境应力剖面分析系统 \\ 详细设计说明书}}
\author{浙江理工大学 \quad 版本 V2.0}
\date{\today}

\begin{document}
\maketitle
\newpage
\tableofcontents
\newpage

% ================================================================
\section{引言}
% ================================================================

\subsection{项目概述}

"船舶平台综合环境应力剖面分析系统"是面向船舶平台设备环境应力监测与数据分析的专业工具。系统基于MATLAB App Designer框架开发，集成数据导入、统计分析、图形化展示等功能，为船舶平台设备的应力状态评估和维护决策提供数据支持。

V1.0系统于2023年11月交付，包含以下核心模块：
\begin{itemize}
    \item \textbf{Tab\_1 振动分析模块}：加速度信号时域/频域分析，功率谱密度计算，规范谱生成
    \item \textbf{Tab\_2 冲击响应分析模块}：冲击响应谱(SRS)计算，多通道包络分析
    \item \textbf{Tab\_3 温湿度分析模块1}：温度/湿度数据前处理、滤波、循环计数
    \item \textbf{Tab\_4 温湿度分析模块2}：温度/湿度综合分析、等效循环计算
    \item \textbf{Tab\_5 倾角分析模块}：三轴倾角时域分析、极值检测
\end{itemize}

V2.0版本在此基础上新增\textbf{Tab\_6 设备应力综合分析模块}，实现对多设备类型、多环境应力参数的自动分类统计与可视化分析。

\subsection{编写目的}

本文档为V2.0新增功能的完整详细设计说明书，旨在：
\begin{enumerate}[label=(\arabic*)]
    \item 明确新增模块的功能需求与设计目标
    \item 详细描述数据流、算法设计和UI交互逻辑
    \item 规定代码集成规范与接口约定
    \item 提供测试方案与验收标准
    \item 作为后续维护和扩展的技术参考
\end{enumerate}

\subsection{读者对象}

\begin{itemize}
    \item \textbf{开发工程师}：理解模块架构和代码实现细节，完成代码集成
    \item \textbf{测试工程师}：依据设计规格编写测试用例，验证功能完整性
    \item \textbf{运维人员}：了解系统结构，进行环境部署和故障排查
    \item \textbf{项目管理人员}：掌握功能边界和交付物清单
\end{itemize}

\subsection{V2.0新增功能需求}

\begin{enumerate}[label=\textbf{FR-\arabic*:}]
    \item \textbf{三列数据导入}：支持.xlsx/.xls/.csv/.txt格式，数据格式为\{设备标签 + 序号, 数值, 单位\}
    \item \textbf{5参数模块}：分别处理蒸汽压力(MPa)、蒸汽流量(m\textsuperscript{3}/h)、蒸汽颗粒物浓度(\%)、盐雾浓度(mg/m\textsuperscript{3})、氧浓度(\%)
    \item \textbf{自动设备分类}：从标签字符串（如"球阀1"）中提取类型名（"球阀"），按类型分组
    \item \textbf{6项统计计算}：最大值(max)、最小值(min)、平均值(mean)、均方根值(rms)、方差(var)、标准差(std)
    \item \textbf{双图可视化}：折线图（时序对比）+ 箱线图（分布分析）
    \item \textbf{可调Y轴范围}：自动适配数据量级，支持手动调节
    \item \textbf{多设备对比}：支持多选设备类型进行对比分析
    \item \textbf{图表导出}：支持PNG和FIG格式导出
    \item \textbf{完整异常处理}：文件校验、数据验证、参数合法性检查、errordlg弹窗提示
\end{enumerate}

\subsection{术语与缩写}

\begin{table}[H]
\centering
\begin{tabular}{c|l}
\toprule
\textbf{术语/缩写} & \textbf{说明} \\
\midrule
Tab\_6 & BoatMain主界面新增的第6个标签页 \\
三列格式 & \{设备标签, 数值, 单位\}的数据结构 \\
设备标签 & 类型名+数字序号，如"球阀1" \\
类型名 & 标签去除末尾数字/分隔符后的前缀，如"球阀" \\
RMS & Root Mean Square，均方根值，有效值 \\
时序测量 & 同一设备在不同时间点的多次测量记录 \\
单位匹配 & 按文件单位列自动筛选对应参数模块的数据 \\
\bottomrule
\end{tabular}
\end{table}

% ================================================================
\section{总体设计}
% ================================================================

\subsection{系统架构}

系统采用MATLAB App Designer单窗口多标签页（TabGroup）架构。V2.0在原有5个标签页基础上，新增第6个\"设备应力分析\"标签页，不影响原有模块功能。

\begin{figure}[H]
\centering
\begin{tikzpicture}[node distance=0.8cm, auto,
    hdr/.style={rectangle, fill=blue!22, text width=18cm, text centered, minimum height=0.9cm, font=\bfseries},
    blk/.style={rectangle, draw, fill=blue!8, text width=3.5cm, text centered, minimum height=0.7cm, font=\normalsize},
    nblk/.style={rectangle, draw, fill=red!12, text width=3cm, text centered, minimum height=0.7cm, font=\bfseries}]

    \node[hdr] (sys) {船舶平台综合环境应力剖面分析系统 V2.0};
    \node[blk, below=0.6cm of sys] (login) {登录认证 (BoatBt.mlapp)};
    \node[blk, below=0.4cm of login] (main) {主界面 (BoatMain.mlapp)};
    \node[blk, below=0.4cm of main] (tg) {TabGroup 标签页容器};

    \node[blk, below left=0.2cm of tg, text width=2.2cm] (t1) {振动分析};
    \node[blk, below=0.1cm of tg, text width=2.2cm] (t2) {冲击响应};
    \node[blk, below right=0.2cm of tg, text width=2.2cm] (t3) {温湿度分析1};
    \node[blk, below=1.2cm of tg, text width=2.2cm] (t4) {温湿度分析2};
    \node[blk, below right=1.2cm of tg, text width=2.2cm] (t5) {倾角分析};
    \node[nblk, below=0.4cm of t4, text width=2.8cm] (t6) {设备应力分析(新增)};
\end{tikzpicture}
\caption{系统架构图}
\end{figure}

\subsection{Tab\_6 模块内部架构}

Tab\_6模块内部按功能分为4个区域：

\begin{figure}[H]
\centering
\begin{tikzpicture}[
    zone/.style={rectangle, draw, dashed, fill=gray!3, minimum width=3.2cm, minimum height=2.2cm, align=center, font=\small},
    title/.style={font=\bfseries}]

    \node[zone, fill=blue!8, minimum width=14.5cm, minimum height=0.6cm] at (0,5.3) {\textbf{标题栏}};

    \node[zone, minimum width=2.5cm, minimum height=4.8cm] at (-5,2.3) {\textbf{控制区}\\[2pt]
    5个模块按钮\\[2pt]
    文件选择\\[2pt]
    设备列表\\[2pt]
    Y轴调节\\[2pt]
    导出/刷新};

    \node[zone, minimum width=5cm, minimum height=2.3cm] at (0.5,3.8) {\textbf{折线图区}\\[2pt]
    设备应力时序对比};

    \node[zone, minimum width=5cm, minimum height=2.3cm] at (0.5,1.3) {\textbf{箱线图区}\\[2pt]
    应力分布分析};

    \node[zone, minimum width=4cm, minimum height=4.8cm] at (5.5,2.3) {\textbf{结果区}\\[2pt]
    统计摘要\\[2pt]
    6项统计表格};
\end{tikzpicture}
\caption{Tab\_6 内部布局分区图}
\end{figure}

\subsection{5个参数模块的独立性设计}

每个参数模块对应一个独立的测试数据文件，模块之间互不依赖：

\begin{table}[H]
\centering
\caption{参数模块定义表}
\begin{tabular}{c|c|c|c|c}
\toprule
\textbf{编号} & \textbf{参数名称} & \textbf{单位} & \textbf{测试文件} & \textbf{按钮控件} \\
\midrule
M1 & 蒸汽压力 & MPa & 蒸汽压力.csv & Button\_18 \\
M2 & 蒸汽流量 & m\textsuperscript{3}/h & 蒸汽流量.csv & Button\_21 \\
M3 & 蒸汽颗粒物浓度 & \% & 蒸汽颗粒物浓度.csv & Button\_22 \\
M4 & 盐雾浓度 & mg/m\textsuperscript{3} & 盐雾浓度.csv & Button\_23 \\
M5 & 氧浓度 & \% & 氧浓度.csv & Button\_24 \\
\bottomrule
\end{tabular}
\end{table}

\textbf{模块切换流程}：用户点击不同模块按钮时，系统自动切换CurParam和CurUnit，重新加载对应文件的数据并执行完整的预处理→分类→统计→绘图流程。

% ================================================================
\section{数据流设计}
% ================================================================

\subsection{数据输入格式规范}

\subsubsection{文件格式要求}

\begin{itemize}
    \item 支持扩展名：.xlsx, .xls, .csv, .txt
    \item 编码：UTF-8（推荐）或GBK
    \item CSV分隔符：逗号
    \item TXT分隔符：制表符(Tab)
    \item 每个参数模块使用独立的文件
    \item 同一文件内所有数据行的单位列应保持一致
\end{itemize}

\subsubsection{三列数据结构}

\begin{table}[H]
\centering
\caption{数据列定义}
\begin{tabular}{c|c|c|c|l}
\toprule
\textbf{列序号} & \textbf{列名} & \textbf{数据类型} & \textbf{是否必填} & \textbf{说明} \\
\midrule
1 & 设备标签 & 字符串 & 必填 & 类型名+序号，如"球阀1" \\
2 & 数值 & 数字 & 必填 & 应力测量值 \\
3 & 单位 & 字符串 & 可选 & 如"MPa"，为空则不过滤 \\
\bottomrule
\end{tabular}
\end{table}

\subsubsection{数据示例}

以蒸汽压力模块为例，文件"蒸汽压力.csv"内容：

\begin{lstlisting}[caption={蒸汽压力测试数据示例}]
设备标签,数值,单位
球阀1,2.5,MPa
球阀1,2.7,MPa
球阀1,2.3,MPa
球阀1,2.9,MPa
球阀2,3.1,MPa
球阀2,2.9,MPa
球阀2,3.0,MPa
阀门1,4.2,MPa
阀门1,4.4,MPa
\end{lstlisting}

\subsubsection{标签命名规则与分组规则}

\textbf{规则很简单}：标签相同的行归为同一组。

例如"球阀1"出现10次（10个时间点的测量值），这10条全部归入"球阀1"组。每个设备标签自己是一组，组与组之间相互独立。

\begin{table}[H]
\centering
\caption{分组示例}
\begin{tabular}{c|c|c}
\toprule
\textbf{数据行} & \textbf{所属组} & \textbf{说明} \\
\midrule
球阀1 的第1$\sim$10行 & 球阀1 & 10条时序记录汇入一组 \\
球阀2 的第1$\sim$10行 & 球阀2 & 另一组 \\
阀门1 的第1$\sim$10行 & 阀门1 & 另一组 \\
\bottomrule
\end{tabular}
\end{table}

\textbf{特殊情况}：标签为空的行自动归为"未命名"组。

\subsubsection{时序数据特性}

同一设备标签在文件中可重复出现任意多次，每次代表该设备在不同时间点的测量值。\textbf{不同设备允许有不同数量的测量记录}，系统自动处理变长数据：

\begin{table}[H]
\centering
\caption{时序数据说明}
\begin{tabular}{c|c|c|c}
\toprule
\textbf{标签} & \textbf{记录数} & \textbf{说明} & \textbf{统计方式} \\
\midrule
球阀1 & 4条 & 4个时间点 & 4个值汇总计算6项统计 \\
球阀2 & 3条 & 3个时间点 & 3个值汇总计算6项统计 \\
阀门1 & 2条 & 2个时间点 & 2个值汇总计算6项统计 \\
\bottomrule
\end{tabular}
\end{table}

MATLAB统计函数（max, min, mean, rms, var, std）原生支持任意长度向量输入，各组数据量不等不影响计算正确性。

\subsection{数据处理流程（完整版）}

\begin{figure}[H]
\centering
\begin{tikzpicture}[node distance=0.65cm, auto,
    start/.style={ellipse, draw, fill=blue!10, minimum width=1.5cm, text centered, font=\small},
    proc/.style={rectangle, draw, fill=green!8, text width=3cm, text centered, minimum height=0.6cm, rounded corners, font=\small},
    check/.style={diamond, draw, fill=yellow!15, text width=2cm, text centered, minimum height=0.8cm, aspect=2, font=\small},
    data/.style={rectangle, draw, fill=orange!10, text width=2cm, text centered, minimum height=0.5cm, font=\small},
    arrow/.style={->, >=stealth, thick},
    line/.style={draw, thick}]

    \node[start] (s) {用户点击模块按钮};

    \node[proc, below of=s] (p1) {检查文件路径\\是否有效};
    \node[check, below of=p1] (c1) {路径非空\\且文件存在?};
    \node[data, right=2cm of c1] (e1) {errordlg弹窗};

    \node[proc, below of=c1] (p2) {readtable读取文件\\自动推断格式};
    \node[check, below of=p2] (c2) {读取成功?};
    \node[data, right=2cm of c2] (e2) {errordlg弹窗};

    \node[proc, below of=c2] (p3) {三列解析:\\col1→tag, col2→value, col3→unit};
    \node[check, below of=p3] (c3) {列数≥2\\且数值列有效?};
    \node[data, right=2cm of c3] (e3) {errordlg弹窗};

    \node[proc, below of=c3] (p4) {按标签前缀分类:\\regexprep提取类型名};
    \node[proc, below of=p4] (p5) {按单位筛选:\\匹配当前模块单位};
    \node[proc, below of=p5] (p6) {6项统计计算};
    \node[proc, below of=p6] (p7) {更新UI:\\表格+折线图+箱线图};
    \node[start, below of=p7] (end) {完成};

    \draw[arrow] (s) -- (p1);
    \draw[arrow] (p1) -- (c1);
    \draw[arrow] (c1) -- node[right]{是} (p2);
    \draw[arrow] (c1) -- node[above]{否} (e1);
    \draw[arrow] (p2) -- (c2);
    \draw[arrow] (c2) -- node[right]{是} (p3);
    \draw[arrow] (c2) -- node[above]{否} (e2);
    \draw[arrow] (p3) -- (c3);
    \draw[arrow] (c3) -- node[right]{是} (p4);
    \draw[arrow] (c3) -- node[above]{否} (e3);
    \draw[arrow] (p4) -- (p5);
    \draw[arrow] (p5) -- (p6);
    \draw[arrow] (p6) -- (p7);
    \draw[arrow] (p7) -- (end);
\end{tikzpicture}
\caption{Tab\_6 数据处理完整流程图（含异常分支）}
\end{figure}

\subsection{自动分类算法（详细）}

\begin{lstlisting}[caption={标签分类核心算法 (MATLAB伪代码)}, label={lst:classify}]
function typeNames = extractTypeNames(tags)
    % 输入: tags - 设备标签字符串cell数组, 如 {'球阀1','球阀2','阀门1'}
    % 输出: typeNames - 类型名cell数组, 如 {'球阀','球阀','阀门'}

    n = length(tags);
    typeNames = cell(n, 1);

    for i = 1:n
        tag = char(string(tags{i}));

        if isempty(tag)
            typeNames{i} = '未分类';          % 空标签处理
            continue;
        end

        % 正则: 去除末尾的数字、连字符、下划线、空格
        cleaned = regexprep(tag, '[\d\-_\s]+$', '');

        if isempty(cleaned)
            typeNames{i} = '数字标签';        % 纯数字标签
        else
            typeNames{i} = cleaned;            % 正常类型名
        end
    end

    % 获取唯一类型并分组
    [uniqueTypes, ~, typeIdx] = unique(typeNames, 'stable');
end
\end{lstlisting}

\textbf{算法复杂度}：$O(n)$，其中$n$为数据总行数。每个标签仅正则处理一次。

\subsection{参数-单位映射机制}

系统通过\textbf{单位模糊匹配}确定数据属于哪个参数模块。匹配规则：数据文件中的单位字符串与模块预定义单位进行双向包含匹配。

\begin{table}[H]
\centering
\caption{参数-单位映射表}
\begin{tabular}{c|c|c|c}
\toprule
\textbf{模块} & \textbf{主单位} & \textbf{兼容单位} & \textbf{匹配逻辑} \\
\midrule
蒸汽压力 & MPa & MPa, MPa(g), Mpa & contains(u, 'MPa') 或 contains('MPa', u) \\
蒸汽流量 & m\textsuperscript{3}/h & m3/h, m\textsuperscript{3}/h & contains(u, 'm3') 或 contains(u, 'm\textsuperscript{3}') \\
蒸汽颗粒物浓度 & \% & \%, pct & contains(u, '\%') 或 contains(u, 'pct') \\
盐雾浓度 & mg/m\textsuperscript{3} & mg/m3, mg/m\textsuperscript{3} & contains(u, 'mg/m') \\
氧浓度 & \% & \%, pct & 同上 \\
\bottomrule
\end{tabular}
\end{table}

% ================================================================
\section{模块详细设计}
% ================================================================

\subsection{UI组件完整清单}

遵循原有BoatMain的命名规范（Label\_57起，Button\_17起，UIAxes\_24起），组件如下：

\begin{longtable}{c|l|l|l|l}
\toprule
\textbf{组件名} & \textbf{MATLAB类} & \textbf{位置[x y w h]} & \textbf{用途} & \textbf{回调} \\
\midrule
\endfirsthead
\toprule
\textbf{组件名} & \textbf{MATLAB类} & \textbf{位置} & \textbf{用途} & \textbf{回调} \\
\midrule
\endhead
\multicolumn{5}{c}{\textbf{标签页容器}} \\
\hline
Tab\_6 & uitab & [1 1 1280 720] & 父容器 & - \\
\hline
\multicolumn{5}{c}{\textbf{标题区域}} \\
\hline
Label\_74 & uilabel & [60 655 450 31] & 模块标题 & - \\
Image\_7 & uiimage & [1 648 43 43] & 图标 & - \\
Label\_76 & uilabel & [60 635 400 20] & 操作提示 & - \\
\hline
\multicolumn{5}{c}{\textbf{文件选择区}} \\
\hline
Button\_17 & uibutton & [42 605 88 24] & 选择数据 & Button\_17Pushed \\
EditField\_17 & uieditfield & [137 606 250 22] & 路径显示 & - \\
\hline
\multicolumn{5}{c}{\textbf{5个参数模块按钮}} \\
\hline
Label\_57 & uilabel & [42 575 200 22] & "参数模块选择" & - \\
Button\_18 & uibutton & [42 546 130 24] & 蒸汽压力 MPa & Button\_18Pushed \\
Button\_21 & uibutton & [180 546 130 24] & 蒸汽流量 m3/h & Button\_21Pushed \\
Button\_22 & uibutton & [42 516 130 24] & 蒸汽颗粒物浓度 & Button\_22Pushed \\
Button\_23 & uibutton & [180 516 130 24] & 盐雾浓度 & Button\_23Pushed \\
Button\_24 & uibutton & [42 486 130 24] & 氧浓度 & Button\_24Pushed \\
\hline
\multicolumn{5}{c}{\textbf{设备列表}} \\
\hline
Label\_58 & uilabel & [42 458 150 22] & "设备类型列表" & - \\
ListBox\_1 & uilistbox & [42 265 150 188] & 设备类型(多选) & ListBox\_1ValueChanged \\
\hline
\multicolumn{5}{c}{\textbf{Y轴调节}} \\
\hline
Label\_73 & uilabel & [42 240 120 22] & "Y轴范围调节" & - \\
EditField\_18 & uieditfield & [100 213 80 22] & Y轴最小值 & EditField\_18ValueChanged \\
EditField\_19 & uieditfield & [100 183 80 22] & Y轴最大值 & EditField\_19ValueChanged \\
\hline
\multicolumn{5}{c}{\textbf{操作按钮}} \\
\hline
Button\_19 & uibutton & [42 145 100 24] & 导出图表 & Button\_19Pushed \\
Button\_20 & uibutton & [150 145 100 24] & 刷新图表 & Button\_20Pushed \\
\hline
\multicolumn{5}{c}{\textbf{绘图区}} \\
\hline
UIAxes\_24 & uiaxes & [230 370 500 260] & 折线图 & - \\
UIAxes\_25 & uiaxes & [230 100 500 260] & 箱线图 & - \\
\hline
\multicolumn{5}{c}{\textbf{结果区}} \\
\hline
UITable\_2 & uitable & [750 240 500 430] & 6项统计表 & - \\
Label\_61 & uilabel & [750 675 300 22] & "统计结果(6项)" & - \\
Label\_62 & uilabel & [750 655 300 22] & 当前模块 & - \\
Label\_63 & uilabel & [750 635 150 22] & 数据条数 & - \\
Label\_64 & uilabel & [920 635 150 22] & 设备类型数 & - \\
\hline
\multicolumn{5}{c}{\textbf{状态栏}} \\
\hline
Label\_66 & uilabel & [42 75 500 22] & 状态信息 & - \\
Label\_75 & uilabel & [42 50 500 22] & 版本信息 & - \\
\bottomrule
\caption{Tab\_6 UI组件完整清单}
\end{longtable}

\subsection{6项统计指标（完整定义）}

新增模块为每个设备类型汇总该类型下所有时序测量值，计算以下6项指标：

\begin{table}[H]
\centering
\caption{6项统计指标完整定义}
\begin{tabular}{c|c|c|c}
\toprule
\textbf{指标} & \textbf{符号} & \textbf{公式} & \textbf{MATLAB} \\
\midrule
最大值 & $x_{\max}$ & $\max(x_1,\ldots,x_n)$ & \texttt{max(v)} \\
最小值 & $x_{\min}$ & $\min(x_1,\ldots,x_n)$ & \texttt{min(v)} \\
平均值 & $\bar{x}$ & $\frac{1}{n}\sum_{i=1}^{n}x_i$ & \texttt{mean(v)} \\
均方根值 & $x_{\text{rms}}$ & $\sqrt{\frac{1}{n}\sum_{i=1}^{n}x_i^2}$ & \texttt{rms(v)} \\
方差 & $s^2$ & $\frac{1}{n-1}\sum_{i=1}^{n}(x_i-\bar{x})^2$ & \texttt{var(v)} \\
标准差 & $s$ & $\sqrt{\frac{1}{n-1}\sum_{i=1}^{n}(x_i-\bar{x})^2}$ & \texttt{std(v)} \\
\bottomrule
\end{tabular}
\end{table}

\textbf{说明}：
\begin{itemize}
    \item \textbf{最大值/最小值}反映该设备类型的应力极值边界
    \item \textbf{均值}反映应力水平的集中趋势
    \item \textbf{均方根值(RMS)}是工程中常用的有效值指标，对较大偏差更敏感
    \item \textbf{方差}和\textbf{标准差}反映应力波动程度，标准差与原始数据同量纲
    \item 使用$n-1$分母（样本方差），保证统计无偏性
\end{itemize}

\subsection{Y轴自动适配算法}

不同模块的数据量级差异很大（如蒸汽压力0$\sim$10 MPa，颗粒物浓度0$\sim$0.3\%），系统在加载数据后自动计算合适的Y轴范围：

\begin{lstlisting}[caption={Y轴自动范围计算}]
% 计算数据范围
allV = app.CurValues(~isnan(app.CurValues));
dataRange = max(allV) - min(allV);
if dataRange == 0, dataRange = 1; end

% Y轴 = [min - 10%margin, max + 10%margin]
ymin = round(min(allV) - 0.1 * dataRange, 3);
ymax = round(max(allV) + 0.1 * dataRange, 3);

app.EditField_18.Value = ymin;  % Y轴最小值
app.EditField_19.Value = ymax;  % Y轴最大值
\end{lstlisting}

用户可通过EditField\_18/19手动调整范围，图表即时刷新。

\subsection{图表自适应规则}

折线图的绘制样式根据每类设备的数据量自动调整，避免标记过密：

\begin{table}[H]
\centering
\caption{折线图自适应规则}
\begin{tabular}{c|c|c|c}
\toprule
\textbf{数据量 n} & \textbf{标记} & \textbf{线宽} & \textbf{适用场景} \\
\midrule
$n \leq 50$ & 圆点 $\bullet$ & 1.5pt & 少量测试数据 \\
$50 < n \leq 200$ & 小点 $\cdot$ & 1.0pt & 中等规模采样 \\
$n > 200$ & 无标记 & 0.8pt & 成百上千条密集数据 \\
\bottomrule
\end{tabular}
\end{table}

箱线图采用全部数据计算，不受数据量影响，天然适合大规模数据分布分析。

\subsection{异常处理设计（完整）}

系统采用\textbf{5层递进式异常处理}策略，每层独立捕获并给出明确提示：

\begin{longtable}{c|l|l|l}
\toprule
\textbf{层} & \textbf{序号} & \textbf{异常场景} & \textbf{处理方式} \\
\midrule
\endfirsthead
\toprule
\textbf{层} & \textbf{序号} & \textbf{异常场景} & \textbf{处理方式} \\
\midrule
\endhead

\multirow{6}{*}{第1层: 文件路径}
& E01 & 文件路径为空 & errordlg + 停止 \\
& E02 & 文件不存在 & errordlg + 停止 \\
& E03 & 文件被其他程序占用 & fopen检测 + errordlg + 停止 \\
\hline

\multirow{4}{*}{第2层: 文件读取}
& E04 & 不支持的文件格式(.pdf/.doc等) & errordlg提示支持的4种格式 + 停止 \\
& E05 & readtable失败(损坏/编码) & errordlg含排查建议 + 停止 \\
& E06 & 文件为空(0行) & errordlg + 停止 \\
\hline

\multirow{6}{*}{第3层: 数据解析}
& E07 & 列数<2 & errordlg提示当前列数 + 正确格式 \\
& E08 & 数值列全部无效 & errordlg + 排查建议 + 停止 \\
& E09 & 部分数值无效 & 逐格转数字, 自动跳过, 状态栏报告数量 \\
& E10 & 标签为空 & 自动替换为"未命名" + 状态栏报告数量 \\
& E11 & 标签含特殊字符 & 正常处理(strtrim后再判断) \\
\hline

\multirow{4}{*}{第4层: 数据筛选}
& E12 & 单位列全为空 & 降级策略: 全部数据纳入分析 \\
& E13 & 无匹配单位的数据 & questdlg询问是否使用全部数据 \\
& E14 & 筛选后无有效数据 & errordlg + 停止 \\
\hline

\multirow{7}{*}{第5层: 统计计算}
& E15 & 某设备数据全空 & 统计值填NaN, 不阻塞其他设备 \\
& E16 & 某设备仅1条(n=1) & 方差/std填NaN; max/min/mean/rms正常 \\
& E17 & RMS计算(无Statistics Toolbox) & fallback: sqrt(mean(v.\^{}2)) \\
& E18 & 设备种类>50 & questdlg询问是否继续 \\
& E19 & Y轴min$\ge$max & 自动重置为数据范围 \\
& E20 & 表格更新失败 & 状态栏警告, 不影响其他输出 \\
\hline

\multirow{3}{*}{第6层: 绘图渲染}
& E21 & 折线图某设备无数据 & 静默跳过该设备 \\
& E22 & 箱线图数据<3个 & 静默跳过该设备 \\
& E23 & 所有数据在Y轴范围外 & 自动扩大范围或静默 \\
\hline

第7层 & E24 & 导出路径无效 & errordlg + 停止导出 \\
\hline
全局 & E99 & 未预期的系统错误 & try-catch兜底 + errordlg + 控制台打印 \\
\bottomrule
\caption{完整异常处理矩阵 (7层24项+1全局)}
\end{longtable}

\textbf{异常处理3条原则}：
\begin{enumerate}
    \item \textbf{不崩溃}：任何异常都不得导致系统崩溃
    \item \textbf{可恢复}：处理后系统回到可操作状态
    \item \textbf{有提示}：用户始终知道发生了什么
\end{enumerate}

% ================================================================
\section{代码结构}
% ================================================================

\subsection{文件组织}

\begin{verbatim}
补充代码/
  BoatMain_Tab6_补充代码.m    # 主文件：5步续写代码

测试数据/
  蒸汽压力.csv               # M1: 球阀/阀门/泵/蒸汽发生器, MPa
  蒸汽流量.csv               # M2: 球阀/阀门/泵/蒸汽发生器, m3/h
  蒸汽颗粒物浓度.csv          # M3: 球阀/阀门/泵/蒸汽发生器, %
  盐雾浓度.csv               # M4: 球阀/阀门/泵/蒸汽发生器, mg/m3
  氧浓度.csv                 # M5: 球阀/阀门/泵/蒸汽发生器, %
\end{verbatim}

\subsection{代码集成5步法}

新增代码完全以"续写"方式注入BoatMain.mlapp，\textbf{不修改任何原有代码}：

\begin{enumerate}[label=\textbf{位置\arabic*:}]

    \item \textbf{properties (Access = public) 末尾}\\
    位置: \texttt{Label\_56} 之后\\
    内容: Tab\_6的所有UI组件声明（Tab\_6, Button\_17--24, EditField\_17--\\
    \quad 19, ListBox\_1, UIAxes\_24--25, UITable\_2, Label\_57--76, Image\_7）

    \item \textbf{properties (Access = private) 末尾}\\
    位置: 私有属性块末尾\\
    内容: DataLoaded, CurParam, CurUnit, CurTags, CurValues,\\
    \quad CurUnits, CategoryNames, CategoryStats

    \item \textbf{createComponents 函数末尾}\\
    位置: \texttt{app.UIFigure.Visible = 'on';} 之前\\
    内容: Tab\_6中所有UI组件的创建和布局代码

    \item \textbf{methods (Access = private) 末尾}\\
    内容: 8个回调函数 + 3个辅助方法\\
    \quad 回调: Button\_17/18/19/20/21/22/23/24Pushed,\\
    \quad\qquad ListBox\_1ValueChanged, EditField\_18/19ValueChanged\\
    \quad 辅助: tab6\_loadAndAnalyze, tab6\_refreshCharts, tab6\_matchType

    \item \textbf{构造函数 (function app = BoatMain)}\\
    位置: \texttt{createComponents(app)} 之前\\
    内容: 状态变量初始化

\end{enumerate}

\subsection{回调函数详细设计}

每个Tab有6个回调函数。以Tab\_6为例，其余Tab\_7$\sim$10结构完全一致，仅控件名后缀不同。

\subsubsection{文件选择 Button\_17Pushed}

\begin{itemize}
    \item \textbf{触发}: 用户点击"选择数据"按钮
    \item \textbf{输入}: 无
    \item \textbf{处理}: 弹出uigetfile对话框（过滤：.xlsx, .xls, .csv, .txt）
    \item \textbf{输出}: 将完整路径写入EditField\_17
    \item \textbf{异常}: 用户取消 → 静默返回
\end{itemize}

\subsubsection{加载分析 Button\_18Pushed → t6\_doAnalyze}

5个Tab的加载按钮全部调用同一个公用函数 \texttt{t6\_doAnalyze}，仅传入的控件句柄不同：

\begin{verbatim}
Button_18 → t6_doAnalyze(app,'蒸汽压力','MPa',
    app.EditField_17, app.ListBox_2, app.UITable_3,
    app.UIAxes_24, app.UIAxes_25,
    app.EditField_18, app.EditField_19,
    app.Label_63...app.Label_68)

Button_24 → t6_doAnalyze(app,'蒸汽流量','m3/h',
    app.EditField_20, app.ListBox_3, ...)
\end{verbatim}

\textbf{t6\_doAnalyze 7层处理流程}：

\begin{enumerate}[label=\textbf{L\arabic*:}, leftmargin=*]
    \item \textbf{文件校验}: 路径空？文件存在？被占用？→ errordlg
    \item \textbf{读取}: 扩展名检查→readtable→空文件检查
    \item \textbf{解析}: c1→标签(空$\to$"未命名"), c2→数值(NaN计数), c3→单位
    \item \textbf{筛选}: 单位匹配→无匹配时questdlg询问→筛选后空?→errordlg
    \item \textbf{分组}: unique(标签)→设备种类>50?→questdlg
    \item \textbf{统计}: 循环计算max/min/mean/rms/var/std, n=1时var/std=NaN, 无rms工具箱时fallback
    \item \textbf{渲染}: 更新ListBox/UITable/Y轴自动/状态Label→调用t6\_doChart绘图
\end{enumerate}

\subsubsection{图表刷新 t6\_doChart}

所有设备列表切换、Y轴参数修改、刷新按钮都调用 \texttt{t6\_doChart}：

\begin{itemize}
    \item \textbf{折线图}: 根据数据量自动选线型（$\le$50条：o-粗线; 50$\sim$200：.-中线; >200：-细线）
    \item \textbf{箱线图}: boxplot，跳过$n<3$的设备
    \item \textbf{匹配}: 直接用 \texttt{strcmp(app.CurTags, 标签名)} 匹配数据
    \item \textbf{Y轴}: 取EditField\_18/19的值，min$\ge$max时自动重置
\end{itemize}

\subsubsection{导出 t6\_export}

\begin{itemize}
    \item uiputfile选择路径和格式(PNG/FIG)
    \item PNG: exportgraphics(ax, path, 'Resolution', 300)
    \item FIG: copyobj→savefig
    \item 异常: errordlg提示
\end{itemize}

\subsection{函数调用关系图}

\begin{verbatim}
用户操作                  回调函数                公用函数
───────                  ────────               ────────
点击"选择数据"  ──→  Button_17Pushed
点击"加载并分析"──→  Button_18Pushed  ──→  t6_doAnalyze()
                                              ├── readtable
                                              ├── unique(分组)
                                              ├── max/min/mean/rms/var/std
                                              ├── 更新UITable_3
                                              ├── 更新Label_63~67
                                              └── t6_doChart()
切换设备选择    ──→  ListBox_2ValueChanged ──→      ├── plot(折线)
修改Y轴范围    ──→  EditField_18ValueChanged ──→    └── boxplot(箱线)
点击"刷新"     ──→  Button_20Pushed ──→
点击"导出"     ──→  Button_19Pushed ──→  t6_export()
\end{verbatim}

\textbf{复用关系}: 5个Tab共享 t6\_doAnalyze、t6\_doChart、t6\_export 三个公用函数，通过传入不同控件句柄区分Tab。

% ================================================================
\section{可调参数设计}
% ================================================================

\subsection{参数列表}

\begin{table}[H]
\centering
\caption{可调参数一览}
\begin{tabular}{c|c|c|c|c}
\toprule
\textbf{参数} & \textbf{UI控件} & \textbf{默认值} & \textbf{范围} & \textbf{生效方式} \\
\midrule
Y轴最小值 & EditField\_18 & 自动计算 & 任意实数 & 修改后即时刷新 \\
Y轴最大值 & EditField\_19 & 自动计算 & 任意实数 & 修改后即时刷新 \\
模块选择 & 5个按钮 & 无默认 & 5选1 & 点击按钮触发重新分析 \\
设备筛选 & ListBox\_1 & 全部选中 & 多选 & 切换后即时刷新 \\
\bottomrule
\end{tabular}
\end{table}

\subsection{自动适配规则}

\begin{enumerate}
    \item 每次加载数据后，系统自动计算 Y = [min - 10\%$\times$range, max + 10\%$\times$range]
    \item 若用户手动修改了Y轴范围且min$\ge$max，系统自动重置为数据范围
    \item 切换设备选择不改变Y轴范围（保持用户设定）
\end{enumerate}

% ================================================================
\section{测试计划}
% ================================================================

\subsection{测试环境}

\begin{itemize}
    \item 操作系统：Windows 10/11
    \item MATLAB版本：R2021a及以上
    \item 必备工具箱：Statistics and Machine Learning Toolbox
\end{itemize}

\subsection{测试数据结构}

每个测试文件包含多个设备标签（如球阀1、球阀2、球阀3、阀门1、阀门2、阀门3、泵1、泵2、泵3），每个设备标签8$\sim$10条时序记录，总计约90条数据。

\subsection{测试用例}

\begin{longtable}{c|l|l|l}
\toprule
\textbf{ID} & \textbf{测试场景} & \textbf{操作步骤} & \textbf{预期结果} \\
\midrule
\endfirsthead
\toprule
\textbf{ID} & \textbf{测试场景} & \textbf{操作步骤} & \textbf{预期结果} \\
\midrule
\endhead

TC01 & 正常加载CSV & 1.点击"蒸汽压力" 2.选择蒸汽压力.csv 3.观察结果 & 表格显示4类设备, 图表正常, Y轴自动适配 \\
\hline
TC02 & 切换模块 & 1.加载蒸汽压力后 2.点击"蒸汽流量" 3.选择蒸汽流量.csv & 数据/图表全部刷新为蒸汽流量数据 \\
\hline
TC03 & 多设备对比 & 1.加载数据后 Ctrl+多选设备 & 图表仅显示选中设备 \\
\hline
TC04 & Y轴调节 & 1.修改EditField\_19值 2.按回车 & 图表Y轴立即更新 \\
\hline
TC05 & Y轴非法值 & 1.设min=100, max=0 & 自动重置为数据范围 \\
\hline
TC06 & 导出PNG & 1.点击"导出图表" 2.选PNG格式 & 生成300dpi PNG文件 \\
\hline
TC07 & 无单位列 & 1.使用无第3列的CSV & 不按单位过滤, 全部数据使用 \\
\hline
TC08 & 部分NaN & 1.第2列含空单元格 & 跳过NaN, 显示警告计数 \\
\hline
TC09 & 文件不存在 & 1.手动输入不存在路径 2.点击模块 & errordlg弹窗 \\
\hline
TC10 & 列数不足 & 1.用单列CSV 2.点击模块 & errordlg弹窗提示列数 \\
\hline
\bottomrule
\caption{测试用例}
\end{longtable}

% ================================================================
\section{附录}
% ================================================================

\subsection{与原系统的兼容性说明}

\begin{itemize}
    \item Tab\_6新增代码与Tab\_1$\sim$Tab\_5完全独立，互不影响
    \item 没有修改任何原有UI组件的Position/回调/属性
    \item 没有修改任何原有数据处理逻辑
    \item 使用独立的properties私有变量，不与原有变量冲突
    \item 回调函数使用tab6\_前缀命名，避免函数名冲突
\end{itemize}

\subsection{版本历史}

\begin{table}[H]
\centering
\begin{tabular}{c|c|c|l}
\toprule
\textbf{版本} & \textbf{日期} & \textbf{作者} & \textbf{变更说明} \\
\midrule
V1.0 & 2023-11-09 & - & 初始版本：5个模块 \\
V2.0 & 2026-07-29 & - & 新增Tab\_6设备应力综合分析模块 \\
\bottomrule
\end{tabular}
\end{table}

\end{document}
